% Bab 1, Topik dari _Aljabar Linear_ Jim Hefferon
%  http://joshua.smcvt.edu/linearalgebra
%  2001-Jun-09
\topic{Menganalisis Jaringan}
\index{jaringan|(}
Diagram di bawah ini memperlihatkan sebagian jaringan listrik sebuah mobil.
Baterai berada di sebelah kiri, digambar sebagai ruas-ruas garis yang bertumpuk. 
Kabel-kabelnya berupa garis, yang digambar lurus dan bersudut siku-siku tajam
agar tampak rapi.
Setiap lampu berupa lingkaran yang mengelilingi sebuah gelung.
\begin{center}
  \includegraphics{gr/mp/ch1.46}
\end{center}
Perancang jaringan seperti itu perlu menjawab pertanyaan seperti:~berapa 
besar arus listrik yang mengalir 
ketika lampu jauh dan lampu rem menyala bersamaan?
Kita akan menggunakan sistem linear untuk menganalisis 
jaringan listrik sederhana. 

Untuk analisis tersebut, kita memerlukan dua fakta tentang listrik 
dan dua fakta tentang jaringan listrik.

Fakta pertama adalah bahwa baterai menyerupai pompa,  
yang memberikan gaya pendorong agar listrik mengalir jika tersedia lintasan.  
Kita mengatakan bahwa baterai memberikan suatu 
\definend{potensial}.\index{potensial, listrik}
% Tentu saja, jaringan ini menjalankan fungsinya ketika, saat listrik  
% mengalir, aliran tersebut melewati sebuah lampu.
Sebagai contoh, ketika pengemudi menginjak rem, sakelar membuat kontak
sehingga membentuk rangkaian pada sisi kiri
diagram, yang mencakup lampu rem. 
Begitu rangkaian terbentuk, gaya dari baterai menghasilkan
aliran melalui rangkaian tersebut, yang disebut \definend{arus}, sehingga lampu menyala.

Fakta kelistrikan kedua adalah bahwa pada beberapa 
jenis komponen jaringan,
besar aliran sebanding dengan gaya yang diberikan oleh baterai.
Artinya, untuk setiap komponen semacam itu terdapat suatu bilangan, yaitu nilai  
\definend{hambatan},\index{hambatan}
sedemikian sehingga potensial sama dengan hasil kali aliran dan hambatan.
Potensial diukur dalam \definend{volt},
laju aliran diukur dalam \definend{ampere},
dan hambatan terhadap aliran diukur dalam \definend{ohm};
satuan-satuan ini didefinisikan sedemikian sehingga
$\mbox{volts}=\mbox{amperes}\cdot\mbox{ohms}$.

Komponen yang memiliki sifat ini,
yaitu kurva respons tegangan-ampere berupa garis yang melalui titik asal,
disebut \definend{resistor}.\index{resistor}
(Tidak setiap komponen memiliki sifat ini.
Sebagai contoh, ketika bola lampu yang ditampilkan di atas 
memanas, nilai hambatannya berubah.)
Sebagai contoh, jika sebuah resistor memiliki hambatan $2$~ohm, 
maka menghubungkannya ke baterai $12$~volt 
menghasilkan arus sebesar $6$~ampere.
Sebaliknya, jika arus listrik sebesar
$2$~ampere mengalir melalui resistor tersebut, harus terdapat 
beda potensial sebesar $4$~volt di antara kedua ujungnya. 
Inilah \definend{penurunan tegangan}\index{penurunan tegangan} pada 
resistor tersebut.
Salah satu cara memandang rangkaian listrik yang kita bahas di sini
adalah bahwa baterai menghasilkan kenaikan tegangan, sedangkan komponen lainnya 
menghasilkan penurunan tegangan.

Fakta tentang jaringan yang kita perlukan adalah  
\textit{Hukum Arus Kirchhoff}, yaitu bahwa pada setiap titik dalam jaringan, arus
     yang masuk sama dengan arus yang keluar, dan
\textit{Hukum Tegangan Kirchhoff}, yaitu bahwa di sepanjang setiap rangkaian, penurunan total sama dengan 
     kenaikan total.\index{Hukum Kirchhoff}\index{jaringan!Hukum Kirchhoff}
% Pada jaringan di atas hanya terdapat satu kenaikan tegangan, yaitu pada baterai,
% tetapi beberapa jaringan memiliki lebih dari satu.

Kita mulai dengan jaringan di bawah ini.
Jaringan tersebut memiliki sebuah baterai yang menyediakan potensial untuk menghasilkan aliran 
dan tiga resistor, yang digambar sebagai garis zig-zag.
Jika komponen-komponen dirangkai secara berurutan seperti di sini,
komponen-komponen itu tersusun \definend{seri}.\index{rangkaian!seri}
\begin{center}
  \includegraphics{gr/mp/ch1.35}
\end{center}
Menurut Hukum Tegangan Kirchhoff, karena kenaikan tegangannya
sebesar $20$~volt, penurunan tegangan total juga harus sebesar $20$~volt.
Karena hambatan dari awal hingga akhir adalah
$10$~ohm (hambatan kabel yang menghubungkan komponen-komponen tersebut dapat diabaikan),
arusnya adalah $(20/10)=2$~ampere. 
Sekarang, menurut Hukum Arus Kirchhoff, arus sebesar $2$~ampere mengalir melalui
setiap resistor.
Dengan demikian, penurunan tegangannya adalah: 
$4$~volt pada resistor $2$~ohm,
$10$~volt pada resistor $5$~ohm, 
dan $6$~volt pada resistor $3$~ohm.

Jaringan sebelumnya cukup sederhana sehingga kita tidak menggunakan sistem linear, tetapi
jaringan berikutnya lebih rumit.
Di sini resistor-resistor tersusun \definend{paralel}.\index{rangkaian!paralel}
% Jaringan ini lebih menyerupai diagram lampu mobil yang ditampilkan sebelumnya.
\begin{center}
  \includegraphics{gr/mp/ch1.36}
\end{center}
Kita mulai dengan memberi label pada cabang-cabang seperti di bawah ini.
Misalkan arus yang melalui cabang kiri dari bagian paralel
adalah $i_1$ dan arus yang melalui
cabang kanan adalah $i_2$, serta
arus yang melalui baterai adalah $i_0$.
% Kita mengikuti Hukum Arus Kirchhoff; sebagai contoh, semua titik pada 
% cabang kanan memiliki arus yang sama, yang kita sebut $i_2$.
Perhatikan bahwa kita tidak perlu mengetahui 
arah aliran yang sebenarnya\Dash jika arus mengalir dalam arah 
yang berlawanan dengan panah kita, 
maka kita akan memperoleh bilangan negatif dalam penyelesaiannya.
\begin{center}
  \includegraphics{gr/mp/ch1.37}
\end{center}
Hukum Arus, yang diterapkan pada titik percabangan di kanan atas, 
memberikan $i_0=i_1+i_2$. 
Jika diterapkan pada titik percabangan di kanan bawah, hukum itu memberikan $i_1+i_2=i_0$.  
Pada rangkaian yang melingkar dari bagian atas baterai, 
turun melalui cabang kiri dari
bagian paralel, lalu kembali ke bagian bawah baterai, kenaikan
tegangannya adalah $20$, sedangkan penurunan tegangannya adalah $i_1\cdot 12$, sehingga
Hukum Tegangan memberikan $12i_1=20$.
Demikian pula, rangkaian dari baterai ke cabang kanan lalu kembali ke
baterai memberikan $8i_2=20$.
Dan, pada rangkaian yang hanya mengelilingi 
cabang kiri dan kanan dari bagian paralel 
(kita memilih arah jarum jam secara sembarang), 
kenaikan tegangannya adalah $0$ dan penurunan tegangannya adalah $8i_2-12i_1$,
sehingga $8i_2-12i_1=0$.
\begin{displaymath}
  \begin{linsys}{3}
   i_0&- &i_1    &-  &i_2   &=  &0 \\
  -i_0&+ &i_1    &+  &i_2   &=  &0  \\
      &  &12i_1  &   &      &=  &20  \\
      &  &       &   &8i_2  &=  &20  \\
      &  &-12i_1 &+  &8i_2  &=  &0  
  \end{linsys}
\end{displaymath}
Penyelesaiannya
adalah $i_0=25/6$, $i_1=5/3$, dan $i_2=5/2$, semuanya dalam ampere.
(Kebetulan, contoh ini menunjukkan bahwa persamaan redundan dapat muncul dalam
praktik.) 

Hukum-hukum Kirchhoff dapat 
menentukan sifat kelistrikan jaringan yang sangat kompleks.
Diagram berikutnya memperlihatkan
lima resistor, dengan nilai dalam ohm, yang dirangkai secara 
\definend{seri-paralel}.\index{rangkaian!seri-paralel} 
\begin{center}
  \includegraphics{gr/mp/ch1.38}
\end{center}
Ini adalah sebuah \definend{jembatan Wheatstone}\index{jembatan Wheatstone} 
(lihat \nearbyexercise{exer:WheatstoneBr}).
Untuk menganalisisnya, kita dapat menempatkan panah-panah dengan cara berikut.
\begin{center}
  \includegraphics{gr/mp/ch1.39}
\end{center}
Hukum Arus Kirchhoff, yang diterapkan pada
simpul atas, simpul kiri, simpul kanan, dan simpul bawah, memberikan
persamaan-persamaan berikut.
\begin{align*}
   i_0     &=  i_1+i_2  \\
   i_1     &=  i_3+i_5  \\
   i_2+i_5 &=  i_4      \\
   i_3+i_4 &=  i_0
\end{align*} 
Hukum Tegangan Kirchhoff,
yang diterapkan pada gelung dalam (gelung $i_0$ ke~$i_1$ ke~$i_3$ ke~$i_0$), 
gelung luar, 
dan gelung atas yang tidak melibatkan baterai, memberikan persamaan-persamaan berikut.  
\begin{align*}
      5i_1+10i_3  &= 10   \\
      2i_2+4i_4   &= 10   \\
      5i_1+50i_5-2i_2  &= 0     
\end{align*} 
Persamaan-persamaan tersebut cukup untuk menentukan penyelesaian 
$i_0=7/3$, $i_1=2/3$, $i_2=5/3$, 
$i_3=2/3$, $i_4=5/3$, dan $i_5=0$. 

Kita dapat memahami banyak jenis jaringan dengan cara ini.
Sebagai contoh, latihan-latihan berikut menganalisis beberapa jaringan jalan.
 
\begin{exercises}
  % \item[] \textit{Banyak sistem untuk soal-soal ini
  %     paling mudah diselesaikan dengan komputer.}
  \item 
    Hitung besar arus dalam ampere pada setiap bagian dari setiap jaringan.
    \begin{exparts}
      \partsitem Ini adalah jaringan sederhana.
          \begin{center}
             \includegraphics{gr/mp/ch1.40}
        \end{center}
      \partsitem Bandingkan jaringan ini dengan kasus paralel yang dibahas di atas.
          \begin{center}
             \includegraphics{gr/mp/ch1.41}
        \end{center}
      \partsitem Ini adalah jaringan yang cukup rumit.
          \begin{center}
           \includegraphics{gr/mp/ch1.42}
        \end{center}
    \end{exparts}
    \begin{answer}
      \begin{exparts}
        \partsitem Hambatan totalnya adalah $7$~ohm.
          Dengan potensial $9$~volt, arusnya akan sebesar $9/7$~ampere.
          Kebetulan, penurunan tegangannya adalah:~$27/7$~volt
          pada resistor $3$~ohm, dan $18/7$~volt pada masing-masing
          dari kedua resistor $2$~ohm.        
        \partsitem Salah satu cara menganalisis jaringan ini adalah dengan memperhatikan bahwa resistor $2$~ohm
          di sebelah kiri mengalami penurunan tegangan sebesar $9$~volt
          (sehingga arus yang melaluinya adalah $9/2$~ampere), dan 
          bagian lain di sebelah kanan juga mengalami penurunan tegangan sebesar 
          $9$~volt, sehingga kita dapat menganalisisnya seperti pada butir sebelumnya.
          Kita juga dapat menggunakan sistem linear.
          \begin{center}
            \includegraphics{gr/mp/ch1.48}
          \end{center}
          Dengan menggunakan variabel-variabel dari diagram, kita memperoleh sistem linear
          \begin{equation*}
            \begin{linsys}{4}
              i_0  &- &i_1  &- &i_2  &  &    &= &0  \\
                   &  &i_1  &+ &i_2  &- &i_3 &= &0  \\
                   &  &2i_1 &  &     &  &    &= &9  \\  
                   &  &     &  &7i_2 &  &    &= &9       
            \end{linsys}
          \end{equation*}
          yang menghasilkan penyelesaian tunggal $i_0=81/14$, $i_1=9/2$, $i_2=9/7$,
          dan $i_3=81/14$.

          Tentu saja, paragraf pertama dan kedua menghasilkan jawaban yang sama.
          Pada dasarnya, dalam paragraf pertama kita menyelesaikan sistem linear 
          dengan metode yang kurang sistematis daripada Metode Gauss, yaitu dengan menentukan beberapa
          variabel lalu melakukan substitusi.  
        \partsitem
          Dengan menggunakan variabel-variabel berikut
          \begin{center}
            \includegraphics{gr/mp/ch1.49}
          \end{center}
          salah satu sistem linear yang cukup untuk menghasilkan penyelesaian tunggal adalah sebagai berikut.
          \begin{equation*}
            \begin{linsys}{7}
              i_0  &- &i_1  &- &i_2  &  &    &  &    & &    & &    &= &0  \\
                   &  &     &  &i_2  &- &i_3 &- &i_4 & &    & &    &= &0  \\
                   &  &     &  &     &  &i_3 &+ &i_4 &-&i_5 & &    &= &0  \\  
                   &  &i_1  &  &     &  &    &  &    &+&i_5 &-&i_6 &= &0  \\  
                   &  &3i_1 &  &     &  &    &  &    & &    & &    &= &9  \\  
                   &  &     &  &3i_2 &  &    &+ &2i_4&+&2i_5& &    &= &9  \\  
                   &  &     &  &3i_2 &+ &9i_3&  &    &+&2i_5& &    &= &9  
            \end{linsys}
          \end{equation*}
          (Tiga persamaan terakhir berasal dari rangkaian yang melibatkan
            $i_0$-$i_1$-$i_6$, 
            rangkaian yang melibatkan $i_0$-$i_2$-$i_4$-$i_5$-$i_6$, 
            dan rangkaian dengan $i_0$-$i_2$-$i_3$-$i_5$-$i_6$.)
           Octave memberikan
            $i_0=4.35616$, $i_1=3.00000$, $i_2=1.35616$,
            $i_3=0.24658$, $i_4=1.10959$, $i_5=1.35616$, $i_6=4.35616$.
      \end{exparts}
    \end{answer}
  \item 
    Pada jaringan pertama yang kita analisis, dengan tiga resistor yang tersusun  
    seri, kita hanya menjumlahkan hambatannya dan mendapati
    bahwa ketiganya bekerja bersama seperti satu resistor $10$~ohm.
    Kita dapat melakukan hal serupa untuk rangkaian paralel. 
    Pada rangkaian kedua yang dianalisis,
    \begin{center}
      \includegraphics{gr/mp/ch1.36}
    \end{center}
    arus listrik yang melalui baterai adalah $25/6$~ampere.
    Dengan demikian, bagian paralel tersebut 
    \definend{ekuivalen}\index{hambatan!ekuivalen}
    dengan satu resistor yang hambatannya 
    $20/(25/6)=4.8$~ohm. 
    \begin{exparts}
      \partsitem Berapakah hambatan ekuivalennya jika resistor
         $12$~ohm diganti dengan resistor $5$~ohm?
      \partsitem Berapakah hambatan ekuivalennya jika kedua resistor masing-masing
         memiliki hambatan $8$~ohm?
      \partsitem Temukan rumus hambatan ekuivalen jika
         kedua resistor paralel masing-masing memiliki hambatan $r_1$~ohm dan $r_2$~ohm.
    \end{exparts}
    \begin{answer}
      \begin{exparts}
        \partsitem 
          Dengan menggunakan variabel-variabel dari analisis sebelumnya,
          \begin{displaymath}
            \begin{linsys}{3}
              i_0&- &i_1    &-  &i_2   &=  &0 \\
             -i_0&+ &i_1    &+  &i_2   &=  &0  \\
                 &  &5i_1   &   &      &=  &20  \\
                 &  &       &   &8i_2  &=  &20  \\
                 &  &-5i_1  &+  &8i_2  &=  &0  
            \end{linsys}
          \end{displaymath}
          Arus yang mengalir di setiap cabang kemudian adalah
          $i_2=20/8=2.5$, $i_1=20/5=4$, dan $i_0=13/2=6.5$, semuanya dalam ampere.
          Dengan demikian, bagian paralel tersebut bekerja seperti satu resistor
          dengan hambatan $20/(13/2)\approx 3.08$~ohm.
        \partsitem 
          Analisis serupa memberikan
          $i_2=i_1=20/8=2.5$ dan $i_0=40/8=5$~ampere.
          Hambatan ekuivalennya adalah $20/5=4$~ohm.
        \partsitem 
          Analisis lain seperti analisis-analisis sebelumnya memberikan 
          $i_2=20/r_2$, $i_1=20/r_1$, 
          dan $i_0=20(r_1+r_2)/(r_1r_2)$, semuanya dalam ampere.
          Jadi, bagian paralel tersebut bekerja seperti satu resistor dengan
          hambatan $20/i_0=r_1r_2/(r_1+r_2)$~ohm.
          (Persamaan ini sering dinyatakan sebagai:~hambatan 
          ekuivalen~$r$ memenuhi $1/r=(1/r_1)+(1/r_2)$.)
      \end{exparts}
    \end{answer}
% Dikomentari karena hambatan bola lampu tidak linear ketika memanas.
%   \item 
%     Untuk contoh dasbor mobil yang membuka Topik ini, tentukan
%     besar arus berikut
%     (anggap bahwa semua hambatan adalah $2$~ohm).
%     \begin{exparts}
%      \partsitem Jika pengemudi menginjak rem, sehingga
%        lampu rem menyala, dan tidak ada rangkaian lain yang tertutup.
%      \partsitem Jika lampu jauh dan lampu rem menyala.
% %      \partsitem Jika lampu jauh, lampu rem, 
% %        dan lampu kabin semuanya menyala.
% %        (Barangkali pengemudi, yang menyadari bahwa lampu kabin menyala sehingga
% %        pintu pasti terbuka, telah menginjak rem!)
%     \end{exparts}
%     \begin{answer}
%       \begin{exparts}
%         \partsitem Rangkaiannya tampak seperti ini.
%           \begin{center}
%             \includegraphics{gr/mp/ch1.52}
%           \end{center}
%         \partsitem Rangkaiannya tampak seperti ini.
%           \begin{center}
%             \includegraphics{gr/mp/ch1.53}
%           \end{center}
%       \end{exparts}
%     \end{answer}
\item \label{exer:WheatstoneBr} 
   Sebuah \definend{jembatan Wheatstone}\index{jembatan Wheatstone}
   digunakan untuk mengukur hambatan. 
   \begin{center}
     \includegraphics{gr/mp/ch1.47}
   \end{center}
   Tunjukkan bahwa pada rangkaian ini, jika arus yang
   mengalir melalui $r_g$ bernilai nol, maka $r_4=r_2r_3/r_1$.
   (Untuk mengoperasikan perangkat ini, tempatkan
   hambatan yang belum diketahui pada $r_4$.
   Pada $r_g$ terdapat alat ukur yang menunjukkan besar arus.
   Kita mengubah-ubah ketiga hambatan $r_1$, $r_2$, dan $r_3$\Dash biasanya
   masing-masing memiliki kenop yang terkalibrasi\Dash hingga 
   alat ukur di tengah menunjukkan $0$.
   Kemudian persamaan tersebut memberikan nilai $r_4$.)
   \begin{answer}
     Hukum Arus Kirchhoff, yang diterapkan pada simpul tempat 
     $r_1$, $r_2$, dan~$r_g$ bertemu, dan juga 
     diterapkan pada simpul tempat $r_3$, $r_4$, dan~$r_g$
     bertemu, memberikan persamaan-persamaan berikut.
     \begin{equation*}
       \begin{linsys}{3}
         i_1 &- &i_2 &- &i_g &= &0  \\
         i_3 &- &i_4 &+ &i_g &= &0  
       \end{linsys}
     \end{equation*}
     Hukum Tegangan Kirchhoff, yang diterapkan pada gelung di kanan atas, 
     dan pada gelung di kanan bawah, memberikan persamaan-persamaan berikut. 
     \begin{equation*}
       \begin{linsys}{3}
         i_3r_3 &- &i_gr_g &- &i_1r_1 &= &0  \\
         i_4r_4 &- &i_2r_2 &+ &i_gr_g &= &0  
       \end{linsys}
     \end{equation*}
     Dengan mengasumsikan bahwa $i_g$ bernilai nol, kita memperoleh $i_1=i_2$, $i_3=i_4$,
     $i_1r_1=i_3r_3$, dan $i_2r_2=i_4r_4$.
     Kemudian, dengan menata ulang persamaan terakhir,
     \begin{equation*}
       r_4=\frac{i_2r_2}{i_4}\cdot\frac{i_3r_3}{i_1r_1}
     \end{equation*}
     dan mencoret faktor-faktor $i$, kita memperoleh kesimpulan yang diinginkan.
   \end{answer}
   %  \item Buktikan Teorema Th\`evenin.
% \item[]\textit{Terdapat jaringan selain jaringan listrik, dan 
%               kita dapat menanyakan seberapa baik hukum-hukum Kirchhoff berlaku padanya.
%               Pertanyaan-pertanyaan selanjutnya membahas perluasan ke 
%               jaringan jalan.}
\item 
    Perhatikan bundaran lalu lintas berikut.
    \begin{center}
      \includegraphics{gr/mp/ch1.43}
     \end{center}
     Berikut adalah volume lalu lintasnya, dalam satuan mobil per sepuluh menit.    
     \begin{center}
       \begin{tabular}{r|ccc}
        \multicolumn{1}{c}{\ } % hilangkan garis vertikal
                            &\textit{North}  
                            &\textit{Pier}  
                            &\textit{Main}  \\
         \cline{2-4}
         \begin{tabular}{@{}r@{}} \textit{masuk} \\ \textit{keluar} \end{tabular}
         &\begin{tabular}{@{}r@{}} $100$ \\ $75$ \end{tabular}
         &\begin{tabular}{@{}r@{}} $150$ \\ $150$ \end{tabular}
         &\begin{tabular}{@{}r@{}} $25$ \\  $50$ \end{tabular}
         % baris:
         % \textit{masuk}      &$100$           &$150$          &$25$      \\
         % \textit{keluar}     &$75$            &$150$          &$50$       
       \end{tabular}
    \end{center}
    Kita dapat menyusun persamaan untuk memodelkan aliran lalu lintas.   
    \begin{exparts}
      \partsitem 
         Sesuaikan Hukum Arus Kirchhoff untuk keadaan ini.
         Apakah asumsi pemodelan ini masuk akal?
      \partsitem 
         Beri label berupa variabel pada ketiga busur di antara jalan-jalan dalam bundaran:~misalkan 
         $i_1$ adalah jumlah mobil  
         yang bergerak dari North Avenue ke Main,
         $i_2$ adalah jumlah mobil di antara
         Main dan Pier,
         dan $i_3$ adalah jumlah mobil di antara Pier dan North.
         Dengan menggunakan hukum yang telah disesuaikan,
         nyatakan sebuah persamaan yang
         menggambarkan aliran lalu lintas pada setiap simpang masuk-keluar.
      \partsitem 
         Selesaikan sistem tersebut.
      \partsitem
         Tafsirkan penyelesaian Anda.
      \partsitem
         Nyatakan kembali Hukum Tegangan untuk keadaan ini.
         Seberapa masuk akalkah hukum tersebut?
    \end{exparts}
    \begin{answer}
      \begin{exparts}
        \partsitem
           Salah satu penyesuaiannya adalah:~pada setiap persimpangan, arus masuk sama dengan 
           arus keluar.
           Hal ini tampaknya masuk akal untuk kasus ini, kecuali jika mobil terhenti di
           suatu persimpangan untuk waktu yang lama.
        \partsitem
           Karena $50$ mobil keluar melalui Main, sedangkan $25$~mobil masuk, 
           $i_1-25=i_2$.
           Demikian pula, keseimbangan masuk/keluar di Pier berarti $i_2=i_3$, dan
           di North diperoleh $i_3+25=i_1$. 
           Kita memperoleh sistem berikut.
           \begin{equation*}
             \begin{linsys}{3}
               i_1  &-  &i_2  &   &     &=  &25  \\
                    &   &i_2  &-  &i_3  &=  &0   \\
              -i_1  &   &     &+  &i_3  &=  &-25
             \end{linsys}
           \end{equation*}
        \partsitem 
           Operasi baris $\rho_1+\rho_2$ dan $\rho_2+\rho_3$ menghasilkan
           kesimpulan bahwa terdapat tak berhingga banyak penyelesaian.
           Dengan $i_3$ sebagai parameter, berikut adalah himpunan penyelesaiannya.
           \begin{equation*}
             \set{\colvec[c]{25+i_3  \\ i_3  \\i_3} \suchthat i_3\in\Re}
           \end{equation*}
           Karena soal dinyatakan dalam jumlah mobil, kita
           mungkin membatasi $i_3$ agar berupa bilangan asli.
        \partsitem
           Jika kita membayangkan bundaran yang awalnya kosong dengan perilaku masukan/keluaran
           seperti yang diberikan, kita dapat menambahkan $i_3$ mobil yang berputar tanpa henti
           untuk memperoleh penyelesaian baru.
        \partsitem
           Salah satu pernyataan ulang yang sesuai adalah:~jumlah mobil yang memasuki 
           bundaran harus sama dengan jumlah mobil yang meninggalkannya.
           Kewajaran pernyataan ini tidak begitu jelas.
           Selama jangka waktu sepuluh menit, kita mungkin mendapati bahwa
           enam mobil lebih banyak masuk daripada keluar, 
           meskipun tabel masuk/keluar dalam pernyataan soal 
           memenuhi sifat ini.
           Bagaimanapun, pernyataan tersebut tidak membantu untuk memperoleh penyelesaian tunggal
           karena untuk itu kita perlu mengetahui jumlah mobil yang berputar
           tanpa henti.
      \end{exparts}
    \end{answer}
  \item 
     Berikut adalah sebuah jaringan jalan. 
          \begin{center}
            \includegraphics{gr/mp/ch1.44}
             \end{center}
         Kita dapat mengamati jumlah mobil per jam yang masuk melalui
         pintu masuk jaringan ini dan keluar melalui pintu keluarnya.
          \begin{center}
            \begin{tabular}{r|ccccc}
               \multicolumn{1}{c}{}  
                 &\textit{Winooski timur}  
                 &\textit{Winooski barat}  
                 &\textit{Willow}  
                 &\textit{Jay} 
                 &\textit{Shelburne} \\ 
                 \cline{2-6}
              \begin{tabular}{@{}r@{}} \textit{masuk} \\ \textit{keluar} \end{tabular} 
              &\begin{tabular}{@{}r@{}} $80$ \\ $30$ \end{tabular} 
              &\begin{tabular}{@{}r@{}} $50$ \\ $5$ \end{tabular} 
              &\begin{tabular}{@{}r@{}} $65$ \\ $70$ \end{tabular} 
              &\begin{tabular}{@{}r@{}} -- \\ $55$ \end{tabular} 
              &\begin{tabular}{@{}r@{}} $40$ \\ $75$ \end{tabular} 
              % baris:
              % \textit{masuk}      &80    &50    &65     &--    &40      \\
              % \textit{keluar}     &30    &5     &70     &55    &75      
            \end{tabular}
          \end{center}
          (Perhatikan bahwa untuk mencapai Jay, sebuah
          mobil harus terlebih dahulu memasuki jaringan melalui jalan lain; karena itu
          tidak ada entri `masuk Jay' dalam tabel.
          Perhatikan pula bahwa dalam jangka waktu yang panjang, 
          jumlah total yang masuk harus kurang lebih sama dengan jumlah total yang 
          keluar; karena itu kedua baris berjumlah $235$~mobil.)
          Setelah berada di dalam jaringan, lalu lintas dapat mengalir dengan berbagai
          cara, mungkin memadati Willow dan membuat Jay 
          hampir kosong, atau mungkin mengalir dengan cara lain.
          Hukum-hukum Kirchhoff memberikan batas-batas bagi kebebasan tersebut.
    \begin{exparts}
       \partsitem Tentukan batasan pada aliran di dalam jaringan 
          jalan ini dengan menetapkan
          sebuah variabel untuk setiap ruas, menyusun persamaan-persamaannya, 
          dan menyelesaikannya.
          Perhatikan bahwa beberapa jalan hanya satu arah.
          (\textit{Petunjuk:}~hal ini tidak akan menghasilkan penyelesaian tunggal, karena lalu lintas
          dapat mengalir melalui jaringan ini dengan berbagai cara;
          Anda akan memperoleh setidaknya satu variabel bebas.)
       \partsitem Misalkan seseorang mengusulkan pekerjaan konstruksi pada
         Winooski Avenue East di antara Willow dan Jay, 
         dan lalu lintas pada ruas itu akan berkurang.
         Berapakah arus lalu lintas paling sedikit yang dapat
         kita izinkan pada ruas tersebut tanpa mengganggu  
         arus per jam yang masuk dan keluar dari jaringan?
    \end{exparts}
    \begin{answer}
      \begin{exparts}
        \partsitem Berikut adalah sebuah variabel untuk setiap ruas yang belum diketahui; setiap ruas yang sudah diketahui
          diberi label besar arusnya.
          \begin{center}
            \includegraphics{gr/mp/ch1.51}          
          \end{center}
          Kita menerapkan prinsip Kirchhoff bahwa arus yang masuk ke persimpangan
          Willow dan Shelburne harus sama dengan arus yang keluar, sehingga diperoleh
          $i_1-i_2=10$.
          Dengan meninjau persimpangan-persimpangan dari kanan ke kiri dan dari atas ke bawah,
          kita memperoleh persamaan-persamaan berikut.
          \begin{equation*}
            \begin{linsys}{7}
              i_1 &- &i_2 &  &    &  &    &  &    &  &    &  &    &= &10  \\
             -i_1 &  &    &+ &i_3 &  &    &  &    &  &    &  &    &= &15   \\
                  &  &i_2 &  &    &+ &i_4 &- &i_5 &  &    &  &    &= &5   \\
                  &  &    &  &-i_3&- &i_4 &  &    &+ &i_6 &  &    &= &-50 \\
                  &  &    &  &    &  &    &  &i_5 &  &    &- &i_7 &= &-10 \\
                  &  &    &  &    &  &    &  &    &  &-i_6&+ &i_7 &= &30 
            \end{linsys}
          \end{equation*}
          Operasi baris $\rho_1+\rho_2$, diikuti oleh $\rho_2+\rho_3$,
          kemudian $\rho_3+\rho_4$, lalu $\rho_4+\rho_5$, dan akhirnya $\rho_5+\rho_6$,
          menghasilkan sistem berikut.
          % runfile("gauss_method.sage")
          % M = matrix(QQ, [[1,-1,0,0,0,0,0], [-1,0,1,0,0,0,0], [0,1,0,1,-1,0,0], [0,0,-1,-1,0,1,0], [0,0,0,0,1,0,-1], [0,0,0,0,0,-1,1]])
          % v = vector(QQ, [10,15,5,-50,-10,30])
          % M_prime = M.augment(v, subdivide=True)
          % gauss_method(M_prime)
          \begin{equation*}
            \begin{linsys}{7}
              i_1 &- &i_2 &  &    &  &    &  &    &  &    &  &    &= &10  \\
                  &  &-i_2&+ &i_3 &  &    &  &    &  &    &  &    &= &25   \\
                  &  &    &  &i_3 &+ &i_4 &- &i_5 &  &    &  &    &= &30  \\
                  &  &    &  &    &  &    &  &-i_5&+ &i_6 &  &    &= &-20  \\
                  &  &    &  &    &  &    &  &    &  &i_6 &- &i_7 &= &-30 \\
                  &  &    &  &    &  &    &  &    &  &    &  &0   &= &0   
            \end{linsys}
          \end{equation*}
          Karena variabel bebasnya adalah $i_4$ dan $i_7$, kita mengambil keduanya sebagai 
          parameter.
          \begin{equation*}
          \begin{split}
            i_6  &=  i_7-30  \\
            i_5  &=  i_6+20=(i_7-30)+20=i_7-10 \\
            i_3  &=  -i_4+i_5+30=-i_4+(i_7-10)+30=-i_4+i_7+20 \\
            i_2  &=  i_3-25=(-i_4+i_7+20)-25=-i_4+i_7-5 \\
            i_1  &=  i_2+10=(-i_4+i_7-5)+10=-i_4+i_7+5
          \end{split}
          \end{equation*}
          Jelas bahwa $i_4$ dan $i_7$ harus tidak negatif, dan
          persamaan pertama menunjukkan bahwa $i_7$ setidaknya bernilai $30$.
          Jika kita mulai dari $i_7$, maka persamaan untuk $i_2$ menunjukkan bahwa
          $0\leq i_4\leq i_7-5$.
        \partsitem Kita tidak dapat mengambil $i_7$ sama dengan nol karena $i_6$ akan
          bernilai negatif (ini berarti mobil bergerak ke arah yang salah di
          Jalan Jay yang satu arah).
          Namun, kita dapat mengambil $i_7$ sekecil $30$, dan dalam hal ini 
          terdapat banyak nilai $i_4$ yang sesuai.
          Sebagai contoh, penyelesaian
          \begin{equation*}
            (i_1,i_2,i_3,i_4,i_5,i_6,i_7)
            =
            (35,25,50,0,20,0,30)
          \end{equation*}
          dihasilkan dengan memilih $i_4=0$.
      \end{exparts}
    \end{answer}
\end{exercises}
\index{jaringan|)}
\endinput
